\documentclass[12pt,a4paper]{article}

\usepackage{michel}
\usepackage{geometry}
\usepackage{amsmath,amssymb,amsthm}
\usepackage{bm}
\geometry{left=2cm, top=1.5cm, right=2cm, bottom=1.5cm}

\usepackage[dcucite,abbr]{harvard}
\harvardparenthesis{none}
\harvardyearparenthesis{round}

\title{Passenger Assignment Model in a Public Transportation Network}
\author{Michel Bierlaire \and Evangelos Paschalidis}
\date{April 23, 2026}

\newtheorem{definition}{Definition}
\newtheorem{assumption}{Assumption}

\begin{document}
\maketitle


\section{Introduction}

This document describes the modeling assumptions, mathematical formulation, and inference components implemented in the software. It is intended as technical documentation for users and developers of the package.

The objective of this software is to infer an aggregate representation of public transportation demand from operational data. The demand is characterized by a flow for each origin--destination (OD) pair and each desired departure time interval.

The software follows four main steps:
\begin{enumerate}
\item Develop a detailed network representation of the public transportation system.
\item Formulate an assignment model that maps an OD matrix to passenger flows on the network.
\item Define measurement equations linking model outputs to observed data, yielding a likelihood function.
\item Estimate demand using ML, MAP, or Bayesian variational inference.
\end{enumerate}

\section{Modeling Elements}

The following sections document the mathematical objects and modeling assumptions used internally by the software.

\subsection{Spatial Representation}

Let $\mathcal{S}$ denote the set of stops in the area of interest. Each stop can act as an origin, a destination, or a transfer location.

\begin{definition}[OD pair]
An origin--destination pair is a tuple
\[
q = (s_o, s_d) \in \mathcal{Q} = \mathcal{S} \times \mathcal{S}.
\]
\end{definition}

\subsection{Temporal Representation}

The analysis period is partitioned into a set $\mathcal{T}$ of departure time intervals. Each interval represents a desired departure time window.

\begin{definition}[Departure time interval]
Each time bin $t\in\mathcal{T}$ corresponds to a closed interval
\[
[a_t,b_t]\subset\mathbb{R},
\qquad a_t<b_t,
\]
where $a_t$ and $b_t$ denote the start and end time (in minutes) of the
desired departure interval.
\end{definition}

For each departure interval $t\in\mathcal{T}$, the desired departure window is the
interval $[a_t,b_t]$. 

\begin{definition}[OD demand]
For each OD pair $q \in \mathcal{Q}$ and departure time interval $t \in \mathcal{T}$, let
\[
f_q^t \ge 0
\]
denote the number of passengers wishing to depart during interval $t$ from origin $s_o$ to destination $s_d$.
\end{definition}

The full demand vector is denoted by
\[
\bm{f} = \{ f_q^t : q \in \mathcal{Q},\; t \in \mathcal{T} \}.
\]

\subsection{Baseline demand and deviation parameterization (Bayesian layer)}

The model assumes that the analyst can provide an \emph{a-priori} demand matrix
$\bm f^{(0)}=\{f_{q}^{t,(0)}\}$, aligned with the OD/time-bin indexing of the scenario.
Rather than placing a prior directly on the demand levels $\bm f$, we parameterize demand
in terms of \emph{log-deviations} from $\bm f^{(0)}$. This improves numerical stability,
facilitates weakly informative priors, and makes the role of prior information explicit.

\paragraph{Modeling choice (smoothly bounded log-deviation).}
For each OD/time-bin cell, introduce an unconstrained latent variable
$\widetilde{z}_q^t\in\mathbb{R}$. To protect the exponentiation and the downstream
assignment from extreme values while retaining a smooth gradient, define the effective
log-deviation as
\[
z_q^t
=
z_{\max}\tanh\!\left(\frac{\widetilde{z}_q^t}{z_{\max}}\right),
\qquad z_{\max}>0.
\]
The demand used by the assignment is then
\[
f_q^t
=
f_{q}^{t,(0)}\exp(z_q^t).
\]
This transformation guarantees $f_q^t\ge 0$ whenever $f_{q}^{t,(0)}\ge 0$ and preserves the baseline
structure up to multiplicative adjustments. It also implies
\[
|z_q^t|\le z_{\max},
\qquad
\exp(-z_{\max})
\le \frac{f_q^t}{f_q^{t,(0)}}
\le \exp(z_{\max})
\]
for positive baseline cells. The default $z_{\max}=6$ therefore permits demand
multipliers ranging approximately from $1/403$ to $403$. In the software configuration,
this smooth bound is exposed under the historical name \texttt{z\_clip}; despite that
name, the implementation uses a smooth hyperbolic-tangent transformation rather than
hard clipping.

\paragraph{Zero baseline cells.}
If $f_{q}^{t,(0)}=0$, then $f_q^t=0$ for all values of $\widetilde{z}_q^t$.
Therefore, OD cells that are zero in the baseline
matrix remain structurally zero during inference.
In order to allow strictly positive deviations from zero, a small but non zero seed value must be provided: $f_{q}^{t,(0)}=\varepsilon>0$.


\paragraph{Prior scale.}
The raw unconstrained variables, rather than the bounded effective deviations, are
assumed independent and normally distributed:
\[
\widetilde{z}_q^t \sim \mathcal{N}(0,\sigma_z^2).
\]
The parameter $\sigma_z>0$ controls the strength of regularization.
Small values of $\sigma_z$ constrain demand to remain close to the baseline, while large values allow
the effective deviations to approach their smooth bounds.
In the current version of the software, $\sigma_z$ is a fixed
hyperparameter.

The distinction between $\widetilde{z}_q^t$ and $z_q^t$ is important: the
Gaussian density is evaluated at the raw variable $\widetilde{z}_q^t$, whereas
the assignment and all reported demand samples use the transformed value $z_q^t$.


\paragraph{Separation of responsibilities.}
The assignment procedure is treated as a deterministic mapping
\[
(\bm f,\theta)\longmapsto \bm x(\bm f;\theta)
\]
and remains agnostic to priors and baseline information. The construction of $\bm f$ from
$\bm f^{(0)}$ and $\widetilde{\bm z}$ is performed in the \emph{Bayesian layer} (implemented using NumPyro and
stochastic variational inference), which: (i) reads the baseline $\bm f^{(0)}$ provided by the user,
(ii) draws the raw deviations $\widetilde{\bm z}$ from their prior distribution,
(iii) applies the smooth bound to obtain $\bm z$, and (iv) computes $\bm f$ using the
expression above before invoking the assignment. This design avoids introducing additional indexing
conventions and keeps the assignment module independent of the inference engine.


\paragraph{Desired departure interval and admissible boarding window.}
For each OD pair $q\in\mathcal{Q}$ and departure time interval $t\in\mathcal{T}$,
we associate the desired departure \emph{interval} $[a_t,b_t]$ and an admissible
schedule-deviation window of half-width $\Delta_{\max}^{\mathrm{access}}>0$
(user-defined, e.g.\ 15 minutes).


\paragraph{Admissible access links.}
An access link is included in the graph \emph{only if} the departure time $\tau$ is not too far
from the desired interval, namely
\[
\tau \in \bigl[a_t-\Delta_{\max}^{\mathrm{access}},\; b_t+\Delta_{\max}^{\mathrm{access}}\bigr],
\]
where $\tau$ denote the actual departure time (in minutes) of a
boarding at the origin stop. As described later, the boarding is
modeled by a link, that is included if the distance from $\tau$ to the
interval $[a_t,b_t]$ does not exceed
$\Delta_{\max}^{\mathrm{access}}$.

\paragraph{Schedule-deviation penalty.}
When the access link is included, its generalized cost is defined by a (possibly asymmetric)
piecewise-linear penalty with a \emph{flat} (zero-penalty) region inside the interval:
\[
c_{\mathrm{access}}(\tau;t)
=
\beta_{\mathrm{early}}\;\max\!\bigl(0,\,a_t-\tau\bigr)
+
\beta_{\mathrm{late}}\;\max\!\bigl(0,\,\tau-b_t\bigr),
\]
where $\beta_{\mathrm{early}}\ge 0$ and $\beta_{\mathrm{late}}\ge 0$ are user-defined coefficients.
Therefore,
\[
c_{\mathrm{access}}(\tau;t)=0
\quad\text{whenever}\quad \tau\in[a_t,b_t],
\]
and penalties apply only to departures strictly earlier than $a_t$ or strictly later than $b_t$.

Note that the thresholds $\Delta_{\max}^{\mathrm{access}}$ and $\Delta_{\max}^{\mathrm{transfer}}$
are fixed configuration parameters; the assignment is differentiable w.r.t.\ costs and OD flows
given a fixed graph.

\paragraph{Implementation note (departure interval as OD/time-bin data).}
Although centroid-in nodes are not time-tagged in the graph, the desired departure
\emph{interval} $[a_t,b_t]$ remains an attribute of the OD/time-bin pair. In computations,
the access-link penalty is evaluated by combining the interval bounds $(a_t,b_t)$ with the
clock time $\tau$ of the \emph{departure event node} reached by that access link. This yields
a zero penalty for departures within $[a_t,b_t]$ and an early/late penalty based on the
nearest bound otherwise, without introducing backward-in-time links in the acyclic
time-expanded network.

\subsection{Network Representation}

The public transportation system is represented as a directed
time-expanded graph. All timetable times (arrival/departure event times) are expressed in minutes on the same time axis as the departure-time intervals.

A key modeling choice is (i) the separation of each stop centroid into two
distinct nodes (an origin \emph{centroid-in} and a destination \emph{centroid-out}),
and (ii) the split of each timetable event into an \emph{arrival} and a
\emph{departure} node. This second split prevents ill-defined instantaneous
moves (e.g., transferring or boarding at the exact same time as arrival) and
ensures that all links respect a strict chronological (or lexicographic) order,
so that the resulting time-expanded network is acyclic.

\begin{definition}[Nodes]
The node set consists of three families of nodes:
\begin{itemize}
  \item \textbf{Centroid-in nodes (non time-tagged, indexed by departure interval):}
      for each stop $s\in\mathcal{S}$ and for each departure interval $t\in\mathcal{T}$,
      a node
      \[
      (s,t)^{\mathrm{in}}
      \]
      used as a demand-injection handle for passengers departing from stop $s$
      with desired departure interval $t$. These nodes do not correspond to a physical
      time instant; they are placed \emph{before all event nodes} in the topological
      order (conceptually at $-\infty$). The interval index $t$ is used only to
      (i) select admissible boarding events through the access-window rule, and
      (ii) compute schedule-deviation penalties.
  \item \textbf{Centroid-out nodes (non time-tagged):} for each stop $s\in\mathcal{S}$,
      a node $s^{\mathrm{out}}$ representing the point where passengers arriving at stop
      $s$ exit the network. These nodes are not associated with a physical time; they
      are placed \emph{after all event nodes} in the topological order (conceptually
      at $+\infty$) so that the global graph remains acyclic.
  \item \textbf{Timetable event nodes (trip-specific, split into arrival and departure):}
      for each stop $s$, each scheduled arrival time $\tau^{\mathrm{arr}}$ and departure time
      $\tau^{\mathrm{dep}}$ associated with a given vehicle run (trip), we create two nodes
      \[
      (s,\tau^{\mathrm{arr}})^{\mathrm{arr}}
      \qquad\text{and}\qquad
      (s,\tau^{\mathrm{dep}})^{\mathrm{dep}}.
      \]
      These nodes are \emph{trip-specific}: if two trips serve the same stop at the same
      clock time, they still correspond to different event nodes.
\end{itemize}
\end{definition}

\begin{definition}[Lines and event-to-line mapping]
Let $\mathcal{L}$ denote the set of public transport lines. Each timetable event node
(arrival or departure) is associated with exactly one operated line, captured by the mapping
\[
\lambda:\mathcal{V}^{\mathrm{event}}\rightarrow\mathcal{L},
\qquad v\mapsto \lambda(v),
\]
where $\mathcal{V}^{\mathrm{event}}$ is the set of timetable event nodes.
In practice, $\lambda(v)$ is derived from the timetable: event nodes created from a given
trip inherit the line identifier of that trip. Therefore, for a given trip at stop $s$,
the corresponding arrival and departure nodes share the same line label.
\end{definition}

Centroid-in nodes are not time-tagged: they are placed before all timetable
event nodes (conceptually at $-\infty$) so that access links can connect to
departure events that occur either before or after the desired departure window
without violating acyclicity. Centroid-out nodes are not time-tagged and
are placed after all event nodes (conceptually at $+\infty$), preserving a global
topological order.

\paragraph{Remark.}
The separation of centroid-in and centroid-out nodes is essential for the
validity of the forward dynamic-programming loading algorithm. Without
this separation, egress links could break the acyclic structure and violate the topological order implied by time.

\begin{definition}[Links]
The link set $\mathcal{A}$ contains the following link types:
\begin{enumerate}
\item \textbf{Access links (time-windowed):}
      for an OD pair $q=(s_o,s_d)$ and departure interval $t$, access links connect the
      origin centroid-in node $(s_o,t)^{\mathrm{in}}$ to \emph{selected}
      \emph{departure} event nodes at the origin stop,
      \[
      (s_o,t)^{\mathrm{in}} \rightarrow (s_o,\tau)^{\mathrm{dep}},
      \]
only for departure times $\tau$ satisfying
\[
\tau \in \bigl[a_t-\Delta_{\max}^{\mathrm{access}},\; b_t+\Delta_{\max}^{\mathrm{access}}\bigr].
\]

\item \textbf{Ride links:}
      for each trip, ride links connect a departure event at the current stop to an
      arrival event at the next stop,
      \[
      (s_i,\tau_i^{\mathrm{dep}})^{\mathrm{dep}} \rightarrow (s_{i+1},\tau_{i+1}^{\mathrm{arr}})^{\mathrm{arr}}.
      \]

\item \textbf{Dwell/continue links (within-trip at a stop):}
      for each trip and each stop event, a dwell link connects the arrival and departure
      nodes at the same stop,
      \[
      (s,\tau^{\mathrm{arr}})^{\mathrm{arr}} \rightarrow (s,\tau^{\mathrm{dep}})^{\mathrm{dep}}.
      \]
      In the implementation, strictly positive dwell time is enforced by regularizing
      timetable records where arrival equals departure (using a small minimum dwell time).

\item \textbf{Transfer links (inter-line, time-windowed):}
      transfer links connect an \emph{arrival} event node to a later \emph{departure} event node
      at the same stop and represent waiting and switching \emph{between different lines}. A transfer link
      \[
      (s,\tau_1)^{\mathrm{arr}} \rightarrow (s,\tau_2)^{\mathrm{dep}},
      \qquad\text{with}\qquad \tau_2>\tau_1
      \]
      is included only if
      \begin{enumerate}
      \item the two event nodes belong to \emph{different lines}, and
      \item the waiting time satisfies $\tau_2-\tau_1 \le \Delta_{\max}^{\mathrm{transfer}}$.
      \end{enumerate}
      No transfer links are created between events of the same line. In particular,
      transfer links are never created for $\tau_2=\tau_1$.

\item \textbf{Egress links (global graph, destination-gated by masking):}
      the global graph contains egress links from every arrival event node at stop $s$
      to the centroid-out node $s^{\mathrm{out}}$:
      \[
      (s,\tau)^{\mathrm{arr}} \rightarrow s^{\mathrm{out}}.
      \]
      For a given OD pair $q=(s_o,s_d)$, only the destination centroid-out node $s_d^{\mathrm{out}}$
      is enabled as a terminal node; all other egress options are disabled by a destination-dependent mask.
\end{enumerate}
\end{definition}

Each link $\ell\in\mathcal{A}$ is associated with:
\begin{itemize}
\item a generalized cost $c_\ell$,
\item a capacity $u_\ell$ for ride links (stored in the data structure, not yet used in costs).
\end{itemize}

\begin{assumption}[Acyclic time-expanded network]
All links in the time-expanded network are consistent with a strict
topological order. All \emph{event-to-event} links strictly increase
clock time. In addition, centroid-in nodes $(s,t)^{\mathrm{in}}$ are
placed \emph{before} all event nodes (conceptually at $-\infty$), and
centroid-out nodes $s^{\mathrm{out}}$ are placed \emph{after} all event
nodes (conceptually at $+\infty$). With this convention, access and
egress links always respect the global ordering.

With the event split, feasible movements follow the pattern
\[
(s,t)^{\mathrm{in}}
\rightarrow (s,\tau)^{\mathrm{dep}}
\rightarrow (s',\tau')^{\mathrm{arr}}
\rightarrow (s',\tau'')^{\mathrm{dep}}
\rightarrow \cdots
\rightarrow (s_d,\bar{\tau})^{\mathrm{arr}}
\rightarrow s_d^{\mathrm{out}}.
\]
Transfer links satisfy $\tau_2>\tau_1$ and dwell links are enforced to be strictly forward in time
through a minimum dwell regularization. Therefore the directed graph is acyclic and admits
a topological ordering, which allows value functions and flows to be computed using dynamic programming.
\end{assumption}

\subsection{Generalized Cost}

Each passenger seeks to minimize a generalized cost that includes:
\begin{itemize}
\item in-vehicle travel time (ride links),
\item waiting and transfer time (transfer links) and dwell time at stops (dwell/continue links),
\item transfer disutility through coefficient $\beta_{\mathrm{transfer}}$,
\item early/late departure penalty relative to the desired departure \emph{interval} (access links; zero penalty within $[a_t,b_t]$).
\end{itemize}

For a path $p$, the generalized cost is
\[
C(p) = \sum_{\ell \in p} c_\ell.
\]

\subsubsection{Link-type-specific cost formulation}

The generalized cost of each link depends on its type. All coefficients
appearing below are fixed user-defined parameters.

\paragraph{Ride links.}
For a ride link $\ell$ representing the movement of a vehicle between two
stops, connecting a departure event to a subsequent arrival event,
\[
(s_i,\tau_i^{\mathrm{dep}})^{\mathrm{dep}} \rightarrow (s_{i+1},\tau_{i+1}^{\mathrm{arr}})^{\mathrm{arr}},
\]
the generalized cost is equal to the in-vehicle travel time:
\[
c_\ell = \tau_{i+1}^{\mathrm{arr}} - \tau_i^{\mathrm{dep}} .
\]

\paragraph{Dwell/continue links.}
For a dwell/continue link at a stop,
\[
(s,\tau^{\mathrm{arr}})^{\mathrm{arr}} \rightarrow (s,\tau^{\mathrm{dep}})^{\mathrm{dep}},
\]
the generalized cost is the dwell time at the stop:
\[
c_\ell = \tau^{\mathrm{dep}} - \tau^{\mathrm{arr}} .
\]

\paragraph{Transfer links (inter-line, bounded waiting).}
Consider a transfer link $\ell:(s,\tau_1)^{\mathrm{arr}}\rightarrow(s,\tau_2)^{\mathrm{dep}}$
with $\tau_2>\tau_1$.
Such a link is included only if (i) the two events correspond to \emph{different lines}
and (ii) the waiting time $\tau_2-\tau_1$ does not exceed the user-defined threshold
$\Delta_{\max}^{\mathrm{transfer}}$.
When included, its generalized cost is
\[
c_\ell = \beta_{\mathrm{transfer}}\,(\tau_2-\tau_1),
\]
where $\beta_{\mathrm{transfer}}>0$ is a user-defined coefficient.

\paragraph{Access links (bounded schedule deviation with a flat region).}
Consider an access link from the centroid-in node $(s_o,t)^{\mathrm{in}}$
to an origin \emph{departure} event node $(s_o,\tau)^{\mathrm{dep}}$.
The link is included only if
\[
\tau \in \bigl[a_t-\Delta_{\max}^{\mathrm{access}},\; b_t+\Delta_{\max}^{\mathrm{access}}\bigr].
\]
When included, its generalized cost is
\[
c_\ell
=
\beta_{\mathrm{early}}\max(0,a_t-\tau)
+
\beta_{\mathrm{late}}\max(0,\tau-b_t),
\]
with user-defined coefficients $\beta_{\mathrm{early}}\ge 0$ and $\beta_{\mathrm{late}}\ge 0$.

\paragraph{Egress links to centroids.}
For an egress link $\ell:(s,\tau)^{\mathrm{arr}}\rightarrow s^{\mathrm{out}}$, we set $c_\ell=0$.

\paragraph{Interpretation of cost coefficients.}
All time quantities in the model are measured in minutes. Travelers may perceive these minutes differently
depending on context. The coefficients $\beta_{\mathrm{transfer}}$, $\beta_{\mathrm{early}}$ and
$\beta_{\mathrm{late}}$ act as perception weights converting physical time into generalized cost.
These coefficients are fixed and specified by the user; they are not estimated within the inference procedure.
Only the OD demand and optionally the dispersion parameter $\theta$ are estimated.

\section{Assignment Model}

\subsection{Input}

The input is the OD vector $\bm{f}$ (scenario order), which is injected through centroid-in nodes.

\subsection{Assignment Procedure}

The classical assignment procedure based on deterministic shortest paths is not differentiable with respect to link costs and is therefore not well
suited for gradient-based calibration methods. We adopt a differentiable assignment model inspired by the logit-based
approach of \citeasnoun{dial1971probabilistic}, adapted to the acyclic time-expanded public transportation network.

\subsection{Differentiable Dial-style assignment}

Let $\bm{f}$ denote the OD demand. For each OD pair $q$ and departure interval $t$, passengers choose routes
according to a logit model defined on the time-expanded network.

\subsubsection{Capacity constraints (current status)}

Although ride links are associated with nominal capacities $u_\ell$ in the data structure, the current
implementation does not incorporate capacity-dependent penalties into generalized costs. All assignment computations
are therefore performed without congestion or overload effects.

Introducing capacity-aware penalties while preserving differentiability (e.g., using smooth overload functions such
as softplus penalties) is a natural future extension, but it is not part of the current version of the software.


\subsubsection{Logit routing model}

Passengers are assumed to choose routes according to a logit model with dispersion parameter $\theta>0$.
For any path $p$,
\[
\Pr(p) \propto \exp\!\left(-\frac{C(p)}{\theta}\right),
\qquad
C(p)=\sum_{\ell\in p}c_\ell.
\]
When $\theta$ is small, flows concentrate on lowest-cost paths.
When $\theta$ is large, flows spread across multiple attractive alternatives.

\paragraph{Computational and behavioral motivation.}
The logit formulation admits an efficient dynamic-programming implementation on the acyclic time-expanded network.
It is differentiable with respect to generalized costs and, consequently, with respect to OD demand.
This property is essential when the assignment model is embedded in gradient-based calibration and variational inference.

\subsubsection{Differentiable dynamic programming formulation}

Fix an OD pair $q=(s_o,s_d)$ and a departure interval $t$. The \emph{terminal} node is the destination centroid-out node
$s_d^{\mathrm{out}}$, and the value function is initialized as
\[
V(s_d^{\mathrm{out}})=0.
\]

For any other node $i$, let $\mathcal{A}^+_q(i)$ denote the set of outgoing links that are enabled for the destination group.
We define
\[
V(i)
=
-\theta\log
\sum_{\ell=(i\to j)\in\mathcal{A}^+_q(i)}
\exp\!\left(-\frac{c_\ell + V(j)}{\theta}\right).
\]
This recursion is evaluated in reverse topological order using a numerically stable \texttt{logsumexp} formulation.

The probability that a passenger located at node $i$ chooses link $\ell=(i\to j)$ is
\[
P_\ell(i)
=
\frac{
\exp\!\left(-\frac{c_\ell + V(j)}{\theta}\right)
}{
  \sum_{(i\to k)\in\mathcal{A}^+_q(i)}
\exp\!\left(-\frac{c_{ik}+V(k)}{\theta}\right)
}.
\]

\subsubsection{Flow propagation}

Fix an OD pair $q=(s_o,s_d)$ and a departure time interval $t\in\mathcal{T}$.
The corresponding demand $f_q^t$ is injected at the origin centroid-in node $(s_o,t)^{\mathrm{in}}$.

Let $y_i$ denote the flow reaching node $i$ during the forward pass. Initialization is
\[
y_{(s_o,t)^{\mathrm{in}}}=f_q^t,
\qquad
y_i=0 \quad \text{for all } i \neq (s_o,t)^{\mathrm{in}}.
\]
Given transition probabilities $P_\ell(i)$ on outgoing links $\ell=(i\to j)$, link flows are
\[
x_\ell = y_i\, P_\ell(i),
\qquad \ell=(i\to j),
\]
and node flows are propagated in topological order by
\[
y_j \leftarrow y_j + x_\ell,
\qquad \ell=(i\to j).
\]

\paragraph{Destination gating.}
The global graph contains egress links for all stops, but for a given OD pair $q=(s_o,s_d)$
the model enables only egress links leading to $s_d^{\mathrm{out}}$ via a destination-dependent mask.

\paragraph{Batched evaluation.}
For computational efficiency, value functions and flow propagation are computed in batches for OD groups sharing the same destination
(and, depending on the implementation, the same departure interval), enabling vectorized evaluation with JAX.

\subsubsection{Opt-out alternative (future extension)}

In a general formulation, an opt-out alternative can be introduced for each OD pair and departure interval. In the first implementation,
opt-out alternatives are not included: all demand is assumed to be assigned to the public transportation network.

\subsection{Implementation notes for fast JAX evaluation}

The time-expanded network is fixed across evaluations. All graph connectivity data are precomputed once and stored as immutable arrays:
a topological ordering of nodes; compressed representations of outgoing adjacency; and, for each link, tail/head indices and link type.

The backward recursion for the value function and forward propagation for flows are implemented using JAX primitives such as
\texttt{logsumexp} and scan-like passes over the topological order. Variable out-degrees are handled through CSR-like adjacency
and/or padded adjacency with masks, enabling compilation with \texttt{jit}.

\subsubsection{Deterministic indexing requirement}
Reproducibility of inference results requires that the indexing of nodes, links, and OD cells remain strictly deterministic across runs.
Graph construction, link ordering, and OD indexing must therefore follow fixed ordering conventions (e.g., sorted iteration and deterministic
array construction). Any change in indexing invalidates stored inference outputs, as the Bayesian layer relies on a fixed mapping between
parameters and link indices.

\section{Measurement Model and Likelihood}

Let $Y$ denote observed data (e.g., passenger counts on selected links, or aggregates derived from them).
Operational measurements are available only for a subset of quantities. Let $\mathcal{M}$ denote the index set of available measurements.
For each $m\in\mathcal{M}$ we denote by $Y_m$ the observed count.

\paragraph{Predicted measurement from the assignment model.}
The assignment model produces link flows $\bm{x}(\bm f;\theta)=\{x_\ell(\bm f;\theta):\ell\in\mathcal{A}\}$.
A deterministic mapping associates each measurement $m$ with a subset of links $\mathcal{A}_m\subseteq\mathcal{A}$,
so that the model-predicted quantity is the sum of the corresponding link flows:
\[
\lambda_m(\bm f;\theta)
=
\sum_{\ell\in\mathcal{A}_m} x_\ell(\bm f;\theta).
\]
In the current version of the software, aggregation weights are implicitly equal to one.


\paragraph{Asymmetric counting device (undercounting).}
The counting device is asymmetric: it is more likely to \emph{miss} passengers than to count passengers that are not present.
We model this by introducing a detection rate $\rho\in(0,1]$ such that the expected observed count is
\[
\mu_m(\bm f;\theta,\rho) = \rho\,\lambda_m(\bm f;\theta).
\]

\paragraph{Negative binomial likelihood.}
Conditionally on $(\bm f,\theta,\rho,r)$, the observations are assumed independent and
\[
Y_m \mid \bm f,\theta,\rho,r \;\sim\; \text{NegBin}\!\bigl(r,\,p_m(\bm f;\theta,\rho,r)\bigr),
\qquad m\in\mathcal{M},
\]
where $r>0$ is a dispersion parameter and the success probability $p_m$ is defined so that
$\mathbb{E}[Y_m]=\mu_m(\bm f;\theta,\rho)$:
\[
p_m(\bm f;\theta,\rho,r)
=
\frac{r}{r+\mu_m(\bm f;\theta,\rho)}
=
\frac{r}{r+\rho\,\lambda_m(\bm f;\theta)}.
\]
Hence, for $y\in\mathbb{N}_0$,
\[
\Pr(Y_m=y\mid \bm f,\theta,\rho,r)
=
\frac{\Gamma(y+r)}{\Gamma(r)\,y!}
\left(1-p_m(\bm f;\theta,\rho,r)\right)^y
\left(p_m(\bm f;\theta,\rho,r)\right)^r.
\]

\paragraph{Fixed measurement parameters.}
In the current version of the software, the detection rate $\rho$ and the dispersion parameter $r$ are treated as fixed, user-specified
hyperparameters and are not estimated within the Bayesian inference procedure.
 Only the OD demand vector and optionally the dispersion
parameter $\theta$ are inferred. Joint estimation of $(\rho,r)$ is a possible future extension.

\paragraph{Likelihood and normalization.}
The likelihood is the product over observed measurements:
\[
\mathcal{L}\bigl(\{Y_m\}_{m\in\mathcal{M}}\mid \bm f,\theta,\rho,r\bigr)
=
\prod_{m\in\mathcal{M}}
\Pr\!\left(Y_m=y_m^{\mathrm{obs}} \mid \bm f,\theta,\rho,r\right).
\]
The corresponding log-likelihood is
\[
\ell(\bm f,\theta,\rho,r)
=
\sum_{m\in\mathcal{M}}
\left[
\log\Gamma\!\left(y_m^{\mathrm{obs}}+r\right)
-\log\Gamma(r)
-\log\!\left(y_m^{\mathrm{obs}}!\right)
+r\log p_m(\bm f;\theta,\rho,r)
+y_m^{\mathrm{obs}}\log\!\left(1-p_m(\bm f;\theta,\rho,r)\right)
\right].
\]
For posterior computations, terms constant with respect to $(\bm f,\theta)$ may be dropped.

\paragraph{Alternative observation models (future extension).}
Other likelihood choices (e.g., Gaussian noise on counts or Poisson models) may be implemented in the future,
but the current baseline implementation uses the negative binomial likelihood described above.

\section{Free and Frozen OD Demand Cells}

An estimation run may hold a selected subset of OD/time-bin cells at predefined
nonnegative values. Let $\mathcal{I}=\{1,\ldots,n_{\mathrm{OD}}\}$ be the full
assignment indexing, and partition it into the free set $\mathcal{U}$ and the
frozen set $\mathcal{F}$, with
$\mathcal{U}\cap\mathcal{F}=\varnothing$ and
$\mathcal{U}\cup\mathcal{F}=\mathcal{I}$. For each $i\in\mathcal{F}$, let
$c_i\geq 0$ denote its predefined flow. The demand reconstruction is
\[
f_i(\bm z_{\mathcal U})=
\begin{cases}
f_i^{(0)}\exp(z_i), & i\in\mathcal{U},\\
c_i, & i\in\mathcal{F}.
\end{cases}
\]
The smooth bounded transformation described previously is applied to the raw
coordinate associated with each $z_i$, $i\in\mathcal U$.

The frozen cells are not represented by dummy parameters. They occur neither in
the optimizer vector nor in the Bayesian guide, prior, gradient, Hessian, or
covariance matrix. The full vector $\bm f$ is reconstructed only at the boundary
of the assignment model, because assignment necessarily consumes all OD cells.
Thus the statistical dimension is exactly
\[
d=|\mathcal U|+\mathbf{1}\{\theta\text{ is estimated}\},
\]
independently of $|\mathcal F|$. A frozen value may be zero or positive. In
contrast, a free cell requires a strictly positive baseline $f_i^{(0)}$, since
the multiplicative parameterization cannot move away from a zero baseline.

Frozen values are supplied by a sparse CSV file with columns
\texttt{origin\_stop\_id}, \texttt{dest\_stop\_id},
\texttt{time\_bin\_id}, and optionally \texttt{fixed\_flow}. A missing or blank
\texttt{fixed\_flow} denotes zero. Keys are validated against the scenario,
duplicates are rejected, and the free/frozen partition is constructed in the
canonical assignment order. The same immutable partition is passed to ML, MAP,
and Bayesian VI. Saved results include the free indices, frozen indices, and
frozen values so that downstream processing can verify the invariant.
The layout also has its own canonical SHA-256 fingerprint, computed from the OD
keys, partition indices, free baselines, and frozen values. This fingerprint is
stored separately from the network/assignment fingerprint: changing only the
frozen-demand file must therefore change the layout identity even though the
network identity remains unchanged. The canonical JSON payload used to compute
the digest is persisted with the result. On loading, the digest is recomputed
from that payload to detect corruption or manual modification. During
post-processing, the layout is rebuilt from the current scenario and
\texttt{fixed\_demand.csv}; its fingerprint must match the result fingerprint
before assignment or report generation proceeds. The network fingerprint and
layout fingerprint are checked independently.

If $|\mathcal U|=0$ and $\theta$ is fixed, then $d=0$. This is treated as a
deterministic evaluation rather than as an optimization or variational-inference
problem. The likelihood is evaluated once for validation, while NumPyro SVI and
the numerical optimizer are bypassed entirely. Returned parameter, gradient,
Hessian, covariance, and posterior-coordinate arrays are correspondingly empty.

\paragraph{Complexity boundary.}
The software exposes a structural complexity profile so that dimensionality can
be audited without relying on machine-dependent timings. For reduced dimension
$d$ and effective low rank $r_q^{\mathrm{eff}}$, the principal statistical
state sizes are
\[
\begin{array}{c|c}
\text{state} & \text{number of scalar entries}\\ \hline
\text{optimizer vector or gradient} & d\\
\text{diagonal Gaussian guide} & 2d\\
\text{low-rank Gaussian guide} & 2d+d r_q^{\mathrm{eff}}\\
\text{full Gaussian guide} & d+d(d+1)/2\\
\text{dense Hessian or covariance} & d^2.
\end{array}
\]
These statistical quantities depend only on the free cells and optional
$\theta$, not on the number of frozen cells.

There is a separate deterministic cost boundary. Each likelihood evaluation
passes a compact demand vector containing the $n_{\mathrm{free}}$ free cells and
the $n_{\mathrm{fixed},+}$ strictly positive frozen cells. Frozen-zero cells are
absent from that vector, from the padded OD group arrays, and from assignment
demand injection. A destination group is also removed when it contains no
remaining cell. Thus the assignment demand-vector size is
$n_{\mathrm{free}}+n_{\mathrm{fixed},+}$ rather than $n_{\mathrm{OD}}$.
The time-expanded graph and other scenario-wide static arrays remain unchanged,
so forward runtime still depends on the graph and on the surviving destination
groups. Positive frozen flows must participate in assignment predictions; when
$\theta$ is estimated, even their route-choice probabilities change with the
estimated parameter. The full $n_{\mathrm{OD}}$ vector is reconstructed only
after estimation for persistence and reporting.
For every VI, ML, or MAP problem, the implementation reports the total, free,
frozen-zero, and frozen-positive OD counts; the compact assignment-vector size;
and the original, surviving, and removed destination-group counts. These
structural metrics accompany elapsed time so that runtime differences can be
interpreted in terms of the work actually performed by the assignment model.

\section{Maximum Likelihood and Maximum A Posteriori Estimation}

The differentiable assignment and measurement models can be used for point estimation
without constructing a full posterior approximation. Let
\[
\bm\eta=
\begin{cases}
(\widetilde{\bm z}_{\mathcal U},\widetilde u), & \text{if $\theta$ is estimated},\\
\widetilde{\bm z}_{\mathcal U}, & \text{if $\theta$ is fixed},
\end{cases}
\]
denote the vector of raw unconstrained parameters. The same smooth transformations
defined above map $\bm\eta$ to the effective demand vector $\bm f$ and, when
applicable, to $\theta$. Consequently, all estimation methods use the same forward
model and negative binomial likelihood.

\subsection{Maximum likelihood estimation}

The maximum likelihood estimator (MLE) is
\[
\widehat{\bm\eta}_{\mathrm{ML}}
=
\arg\max_{\bm\eta}
\ell(\bm\eta),
\]
where
\[
\ell(\bm\eta)
=
\log\mathcal L\!\left(
Y\mid \bm f(\widetilde{\bm z}),\theta(\widetilde u),\rho,r
\right),
\]
with the dependence on $\widetilde u$ omitted when $\theta$ is fixed. The MLE
uses no prior information. The baseline demand $\bm f^{(0)}$ still appears in
the parameterization $\bm f=\bm f^{(0)}\exp(\bm z)$, but it is not a probabilistic
prior in the ML objective.

After optimization, the raw point estimate is transformed using exactly the same
smooth mappings as in Bayesian inference:
\[
\widehat{z}_q^t
=z_{\max}\tanh(\widehat{\widetilde z}_q^t/z_{\max}),
\qquad
\widehat f_q^t=f_q^{t,(0)}\exp(\widehat z_q^t),
\]
and, when estimated,
\[
\widehat u
=u_{\max}\tanh(\widehat{\widetilde u}/u_{\max}),
\qquad
\widehat\theta=\exp(\widehat u).
\]

\subsection{Maximum a posteriori estimation}

Maximum a posteriori (MAP) estimation uses the same likelihood and priors as the
Bayesian model, but returns only the posterior mode:
\[
\widehat{\bm\eta}_{\mathrm{MAP}}
=
\arg\max_{\bm\eta}
\left\{
\ell(\bm\eta)+\log p(\bm\eta)
\right\}.
\]
Here $p(\bm\eta)$ contains the Gaussian prior on $\widetilde{\bm z}$ and,
when $\theta$ is estimated, the Gaussian prior on $\widetilde u$. The priors
are evaluated on the raw variables, not on their smoothly bounded transforms.

The software implements ML and MAP through the common penalized objective
\[
Q_w(\bm\eta)
=
-\ell(\bm\eta)-w\log p(\bm\eta),
\qquad w\ge 0.
\]
The setting $w=0$ gives pure maximum likelihood and $w=1$ gives MAP under the
same prior as the Bayesian model. Intermediate or larger values provide regularized
point estimates, but do not in general correspond to the posterior of the stated
Bayesian model.

\subsection{Numerical optimization and local uncertainty}

The current point-estimation implementation minimizes $Q_w$ with gradient-based
SciPy optimizers such as BFGS or L-BFGS-B. Gradients are supplied by JAX automatic
differentiation. Because the objective need not be globally convex, convergence
to the global optimum is not guaranteed; initialization, stopping tolerances, and
local curvature should therefore be inspected.

When requested, the software computes the Hessian of $Q_w$ at the optimum and
uses its inverse as a local covariance approximation. For ML this is an asymptotic
likelihood-based approximation; for MAP it is a Laplace-style local approximation
to posterior curvature. These standard errors are not equivalent to a complete
posterior distribution and may be unreliable for weakly identified, asymmetric,
bounded, or multimodal problems.

\section{Bayesian Inference via Variational Inference}

Let $\widetilde{\bm z}_{\mathcal U}$ denote the vector of raw OD-deviation variables for free cells and let
$\widetilde u$ denote the raw variable associated with the dispersion parameter
of the route choice model. The effective parameters supplied to the forward model
are deterministic transformations of these raw variables.
The posterior distribution is proportional to
\[
p(\widetilde{\bm z}_{\mathcal U},\widetilde u \mid Y)
\propto
\mathcal{L}\!\left(Y \mid \bm f(\widetilde{\bm z}_{\mathcal U}),
\theta(\widetilde u)\right)\,
p(\widetilde{\bm z}_{\mathcal U})\,p(\widetilde u).
\]
When $\theta$ is fixed, $\widetilde u$ and its prior are omitted.

Exact Bayesian inference is computationally intractable for high-dimensional OD vectors.
The current implementation therefore uses \emph{stochastic variational inference} (SVI) as implemented in NumPyro.

\paragraph{Variational approximation.}
The posterior is approximated by a parametric distribution
$q_\phi(\widetilde{\bm z}_{\mathcal U},\widetilde u)$ in the raw unconstrained parameter space.
The parameters $\phi$ are chosen to maximize the \emph{Evidence Lower Bound} (ELBO):
\[
\mathrm{ELBO}(\phi)
=
\mathbb{E}_{q_\phi}
\left[
\log p(Y,\widetilde{\bm z}_{\mathcal U},\widetilde u)
-
\log q_\phi(\widetilde{\bm z}_{\mathcal U},\widetilde u)
\right].
\]
Maximizing the ELBO is equivalent to minimizing $\mathrm{KL}(q_\phi \,\|\, p)$ between the variational approximation and the true posterior.

\paragraph{Guide family.}
The baseline implementation uses Gaussian variational families in an unconstrained parameter space
(e.g., diagonal or low-rank approximations), and optimization is performed using stochastic gradients (Adam).
For a requested low-rank value $r_q$, the implemented effective rank is
\[
r_q^{\mathrm{eff}}=\min(r_q,d),
\]
where $d=|\mathcal U|+\mathbf{1}\{\theta\text{ is estimated}\}$ is the reduced
statistical dimension. This prevents a guide configured for the original model
from retaining unnecessary covariance factors after many OD cells are frozen.
The effective rank, rather than the oversized requested value, is recorded in
the inference diagnostics. When $d=0$, no guide is constructed.

\paragraph{Prior on $\theta$.}
The dispersion parameter $\theta$ is strictly positive. When it is estimated,
the implementation introduces an unconstrained raw variable $\widetilde u\in\mathbb R$
and defines
\[
u
=
u_{\max}\tanh\!\left(\frac{\widetilde u}{u_{\max}}\right),
\qquad
\theta=\exp(u),
\qquad u_{\max}>0.
\]
Thus $\theta$ is positive and smoothly bounded between
$\exp(-u_{\max})$ and $\exp(u_{\max})$. The configuration currently exposes
$u_{\max}$ under the historical name \texttt{u\_clip}, although no hard clipping
is performed.

A Gaussian prior is placed on the raw variable,
\[
\widetilde u\sim\mathcal N(\mu_{\widetilde u},\sigma_u^2).
\]
If the requested center in effective log-dispersion space is $\mu_u$---for example,
$\mu_u=\log(\theta^{(0)})$ for a baseline value $\theta^{(0)}$---the raw prior
center is chosen by the inverse transformation,
\[
\mu_{\widetilde u}
=
u_{\max}\operatorname{arctanh}\!\left(\frac{\mu_u}{u_{\max}}\right).
\]
This guarantees that transforming the raw prior center gives exactly the requested
effective center. It requires $|\mu_u|<u_{\max}$.

\paragraph{Posterior transformation and reporting.}
Posterior samples produced by the variational guide are samples of the raw variables
$\widetilde{\bm z}$ and, when applicable, $\widetilde u$. Before reporting demand or
$\theta$, the software applies the same smooth transformations used during likelihood
evaluation:
\[
z_q^t=z_{\max}\tanh(\widetilde z_q^t/z_{\max}),
\qquad
f_q^t=f_q^{t,(0)}\exp(z_q^t),
\]
and
\[
u=u_{\max}\tanh(\widetilde u/u_{\max}),
\qquad \theta=\exp(u).
\]
Consequently, reported posterior summaries correspond exactly to the parameter values
used by the assignment and likelihood, rather than to untransformed raw samples.

\paragraph{Alternative inference engines.}
Markov chain Monte Carlo methods could be used in principle, but they are computationally expensive in high-dimensional OD settings.
Support for alternative probabilistic programming frameworks (e.g., PyMC) may be considered in the future; the present baseline is NumPyro+SVI.

\section{Comparison of the Estimation Methods}

The three estimation approaches share the assignment model, measurement equation,
likelihood, raw parameter space, and smooth parameter transformations. They differ
in the statistical quantity they target.

\paragraph{ML versus MAP.}
The MLE maximizes the likelihood alone, whereas MAP balances the likelihood with
the specified prior. They therefore need not coincide. MAP is typically pulled
toward the prior center and can be more stable in weakly identified problems.
As the amount of informative data increases and the likelihood dominates the prior,
the ML and MAP estimates will often become close, subject to regularity and
identifiability conditions.

\paragraph{MAP versus Bayesian inference.}
If MAP and Bayesian inference use exactly the same likelihood, priors, and
parameterization, the MAP estimate is the mode of the exact Bayesian posterior.
It is not generally equal to the posterior mean or median. Equality between the
mode and mean occurs only in special cases, such as an exactly symmetric unimodal
posterior. The difference may be substantial for skewed, weakly identified, bounded,
or multimodal posteriors.

\paragraph{Variational approximation.}
The implemented Bayesian procedure reports means and other summaries computed from
the variational approximation $q_\phi$, not from the exact posterior. Differences
between MAP and VI results can therefore arise from three distinct sources:
\begin{enumerate}
\item the intrinsic difference between a posterior mode and a posterior mean;
\item approximation error introduced by the selected variational guide; and
\item optimization and finite posterior-sampling error in the VI procedure.
\end{enumerate}
The posterior draws used to compute VI summaries introduce Monte Carlo variability,
but this is not the only, and often not the dominant, reason for differences between
the estimates.

\paragraph{When estimates should be close.}
Under standard regularity conditions, with a large informative sample and a
well-identified model, the posterior is often approximately Gaussian and concentrated.
In that regime the MLE, MAP estimate, posterior mode, and posterior mean may be close.
This is an asymptotic property rather than an equality guaranteed by the software.

\paragraph{Recommended comparisons.}
For validating the implementation, pure ML should be compared across optimizers
using $w=0$, and MAP should be compared with the mode of the same posterior using
$w=1$. Comparing an ML or MAP point estimate directly with the reported VI posterior
mean is useful diagnostically, but disagreement should not automatically be interpreted
as a software error. In every comparison, the effective transformed values
$(\bm f,\theta)$ should be compared rather than the raw variables
$(\widetilde{\bm z},\widetilde u)$.

\section{Discussion}

This software framework combines timetable information and operational counts to infer OD demand and quantify uncertainty,
while using a differentiable assignment model compatible with ML, MAP, and
gradient-based Bayesian inference. ML and MAP provide computationally convenient
point estimates and local curvature diagnostics, whereas Bayesian VI additionally
provides an approximate distribution that represents parameter uncertainty.

Planned extensions of the software include incorporating capacity effects in a
differentiable manner, estimating additional measurement parameters
(e.g., $\rho$ and $r$), supporting weighted measurement aggregations,
and adding opt-out alternatives.


\bibliographystyle{dcu}
\bibliography{transpor}

\end{document}
